\section{Introduction}\label{sec:introduction}

A smooth cubic fourfold \(X\subset\mathbf P^5_{\C}\) is rational if its field of rational functions is
a purely transcendental extension of \(\C\) of transcendence degree four. The connection
between this birational question and K3 surfaces appears both in the
middle cohomology and in the derived category. We prove that both
connections can hold for the same irrational cubic.

The geometric connection grew out of explicit rationality constructions.
Cubics containing quartic rational scrolls were studied by Fano
\cite[p.~72]{Fano43}; Tregub gave constructions using two disjoint
planes, quintic del Pezzo surfaces, and quartic scrolls
\cite[Sections~2--4]{Tregub84}.
Beauville and Donagi proved that the Fano variety \(F(X)\), which
parametrizes lines on \(X\), is irreducible holomorphic symplectic.
For sufficiently general Pfaffian cubics they identified \(F(X)\) with
\(S^{[2]}\), the variety of length-two subschemes of a K3 surface \(S\),
and proved that the cubic is rational
\cite[Propositions~1--2 and~5]{BeauvilleDonagi85}.
Hassett constructed further rational cubics containing a plane; in those
parametrizations a family of exceptional lines is parametrized by a
surface birational to a K3 surface
\cite[Theorem~4.2 and Proposition~5.1]{Hassett99}.

The Kuznetsov component of \(X\) is the full triangulated subcategory
\[
 \Ku(X)=\bigl\{F\in D^b(\operatorname{Coh}X):
 \Hom(\mathcal O_X(i),F[k])=0
 \text{ for }0\le i\le2\text{ and }k\in\Z\bigr\}.
\]
It gives the semiorthogonal decomposition
\[
 D^b(\operatorname{Coh}X)
 =\langle\Ku(X),\mathcal O_X,\mathcal O_X(1),\mathcal O_X(2)\rangle.
\]
Kuznetsov's rationality conjecture asserts that \(X\) is rational if and only
if \(\Ku(X)\) is equivalent to the derived category of a projective K3
surface \cite[Conjecture~1.1]{Kuznetsov10}. Here rationality means
birationality over \(\C\) with \(\mathbf P^4\), and the K3 equivalence in
this statement means an exact \(\C\)-linear equivalence with the ordinary category
\(D^b(\operatorname{Coh}S)\). The target has no Brauer twisting.

Let \(\mathcal C\) denote the moduli space of smooth complex cubic
fourfolds. Write \(h=c_1(\mathcal O_X(1))\). A \emph{labelling of
discriminant \(d\)} is a saturated positive-definite rank-two sublattice
\[
 K\subset H^4(X,\Z)\cap H^{2,2}(X)
\]
containing \(h^2\), with \(\det K=d\). The corresponding Hassett locus is
denoted by \(\mathcal C_d\). We call an integer \(d\) \emph{admissible} if
\begin{equation}\label{cub:admissible}
 d>6,\qquad d\equiv0,2\pmod6,\qquad
 4\nmid d,\quad 9\nmid d,\quad
 p\nmid d\ \text{for every odd prime }p\equiv2\pmod3.
\end{equation}
Hassett's period and lattice theorems make \(\mathcal C_d\) a nonempty
irreducible divisor for these discriminants and give its associated
polarized K3 surfaces \cite[Theorems~1.0.1 and~1.0.2]{Hassett00}.
Here an associated K3 surface has a primitive polarization \(L\) of
degree \(d\) and an integral Hodge isometry
\(K^\perp\simeq L^\perp(-1)\), where the surface cup pairing is reversed.

Kuznetsov verified the categorical prediction for smooth Pfaffian cubics
\cite[Theorem~3.1]{Kuznetsov10}. Addington and Thomas proved that the
Hodge-theoretic and categorical K3 associations agree on a dense open
subset of each admissible divisor \cite[Theorem~1.1]{AddingtonThomas14}.
The construction of stability conditions by Bayer, Lahoz, Macr\`i and
Stellari \cite[Theorem~1.2]{BLMS23}, together with the theory in families
of Bayer, Lahoz, Macr\`i, Nuer, Perry and Stellari, gives an ordinary K3
derived category for every cubic having an associated K3 surface
\cite[Corollary~29.7]{BLMPNS21}. These results supply the categorical
assertion below; the irrationality assertion is proved here.

Our main result disproves the categorical-to-rational implication of
Kuznetsov's conjecture on every sufficiently large admissible divisor.

\begin{theorem}[Irrationality on large admissible Hassett divisors]\label{thm:main}
There is an integer \(d_0\) such that, for every admissible \(d>d_0\), a very
general cubic fourfold \(X\) in \(\mathcal C_d\) is irrational
and admits an exact \(\C\)-linear equivalence
\[
 \Ku(X)\simeq D^b(\operatorname{Coh}S)
\]
for a smooth projective K3 surface \(S\).
\end{theorem}

The threshold \(d_0\) is ineffective. Some Hassett divisors consist
entirely of rational cubics: Russo and Staglian\`o proved rationality throughout
\(\mathcal C_{26}\) and \(\mathcal C_{38}\) using congruences of
five-secant conics \cite[Theorems~4 and~7]{RussoStagliano19}, and
throughout \(\mathcal C_{42}\) using trisecant flops
\cite[Theorem~5.12]{RussoStagliano23}.
Here ``very general'' means outside a
countable union of proper closed algebraic subsets separately in each
divisor \(\mathcal C_d\).
This qualification matters: Yang and Yu construct rational subloci of
codimension two in the full moduli space inside every Hassett divisor
\cite[Theorem~9]{YangYu20}.

\subsection{The two evaluations of one divisorial invariant}

Write \(T_X\) for the integral lattice orthogonal to the integral
Hodge classes in \(H^4(X,\Z)\), with its cup-product pairing, and put
\(d=|\det T_X|\). The irrationality argument concerns cubics with
\(\operatorname{rank}T_X=21\), \(d>2\), and
\(\operatorname{End}_{\mathrm{Hdg}}(T_X\otimes\Q)=\Q\).
Suppose such an \(X\) were rational, and choose a birational map
\(f:\mathbf P^4\dashrightarrow X\).
We use one fixed collection \(\mathcal S\) of
Picard-number-one K3 surfaces: those whose entire integral
transcendental Hodge lattice, after shifting Hodge types by \((1,1)\)
and reversing the pairing, is isometric to \(T_X\).
The lattice hypotheses imply that these surfaces have primitive ample
degree \(d\) and trivial automorphism group
(Lemma~\ref{gw:k3-center}). We do not assume that this collection is
nonempty: the hypothetical birational map will force a member of it.

The maximal rationally connected (MRC) quotient of a smooth projective
variety is the base of its birationally determined rational fibration with
rationally connected general fibres and non-uniruled base.
For an integral variety \(E\), let \(u(E)\) be one if the MRC quotient
of a smooth projective model of its function field is birational to a
member of \(\mathcal S\), and zero otherwise. Section~\ref{pre:section} defines \(c_u(f)\) by counting the
prime divisors acquired by the target with positive sign and those
lost from the source with negative sign, each weighted by \(u\).
The count is finite and additive under composition. We evaluate this
same invariant in two ways.

Smooth weak factorization expresses \(f\) as blowups and blowdowns
along smooth centers. Only surface centers can have nonzero test.
Classically, a surface blowup adds its whole transcendental lattice,
shifted to weight four with reversed pairing, as an integral orthogonal
summand. The difficulty is to recognize that same summand on blowdown.
Section~\ref{gw:section} uses genus-zero Gromov--Witten pairings with
two variable transcendental inputs and all remaining insertions rational
of diagonal Hodge type. For an intermediate fourfold \(M\),
define \(T_M\) as for \(X\). Each such pairing \(B(a,b)\) determines an
operator \(A_B\) by
\(B(a,b)=q_M(A_Ba,b)\), where \(q_M\) is the cup form on
\(T_M\otimes\Q\). The algebra generated by these operators has a
separate scalar factor on each whole surface contribution. Its
projectors force a later blowdown to remove a whole block. Intersecting
that rational block with the ambient integral lattice recovers its
entire integral lattice, so cancellation preserves the exact test
class. Every nonzero block has a nonzero \((3,1)\) part, so
\(h^{3,1}(X)=1\) leaves one block in the final transcendental lattice.
Rank \(21\) makes its surface center birational to a
Picard-number-one K3 surface, and its discriminant fixes the degree
\(d\). Thus one more member of the class \(\mathcal S\) is added than
removed, and Proposition~\ref{gw:net-center} gives \(c_u(f)=1\).

The Sarkisov factorization passes through Mori fibre spaces: contractions
\(V\to B\) with terminal \(\Q\)-factorial projective fourfold \(V\),
\(\dim B<4\), relative Picard number one, and relatively ample
\(-K_V\). These are the nodes of the factorization. An elementary
birational transformation between nodes is a link; its endpoint bases
map to a common base \(Q\). The bases can change along the path.
Sections~\ref{bd:section}--\ref{sf:section} construct integer-valued
functions \(v,s\) on the nodes and prove, for sufficiently large \(d\),
\[
 c_u(\text{link})=(v+s)(V'/B')-(v+s)(V/B).
\]
The value \(v(V/B)\) is the finite cancelled difference between the
tests of vertical divisors on \(V\) and divisors on \(B\).
The value \(s\) comes from the remaining generic-surface calculation.

At a conic-bundle node, the base field \(L=\C(B)\) has transcendence
degree three. For a surface subfield \(k\subset L\) arising from a link,
\(L\) is the function field of a genus-zero curve over \(k\). The surface
calculation
then gives a candidate for \(s\), using the test on finite residue
fields and covers determined by the conic's Brauer ramification.
The crucial point is that the candidates agree for every such
presentation with the ramification bound supplied by the links.
A second Sarkisov factorization, now on the threefold bases, compares
these presentations. Bounded diagrams with their branch loci recorded
give genus bounds incompatible with pencils on Picard-number-one K3
surfaces of sufficiently large degree. This proves independence with
one bound for all test subcollections and factorizations. At nodes over
a point both corrections vanish; the displayed differences therefore
telescope to \(c_u(f)=0\) by Theorem~\ref{sf:vanishing}.

The contradiction proves the irrationality criterion with one threshold
for all the lattice data in its statement. Section~\ref{cub:section} then
applies the standard period and categorical results: for each admissible
\(d\), a very general point of \(\mathcal C_d\) has the required lattice
data and an exact \(\C\)-linear equivalence of its Kuznetsov component
with the ordinary derived category of a smooth projective K3 surface.
The criterion therefore applies in every admissible divisor above the
same threshold. Figure~\ref{fig:two-evaluations} summarizes the
contradictory evaluations of the hypothetical map.

\begin{figure}[ht]
\centering
\begin{tikzcd}[column sep=large,row sep=large]
 & \mathbf P^4\dashrightarrow X
   \arrow[dl,"\text{smooth factorization}" description]
   \arrow[dr,"\text{Mori factorization}" description] & \\
 \text{whole integral surface blocks}
   \arrow[d,"\text{Proposition~\ref{gw:net-center}}" description]
 && \text{endpoint corrections }v+s
   \arrow[d,"\text{Theorem~\ref{sf:vanishing}}" description] \\
 c_u(f)=1 && c_u(f)=0
\end{tikzcd}
\caption{Assuming rationality gives a birational map
\(f:\mathbf P^4\dashrightarrow X\). Its two evaluations use the same
collection of K3 surfaces.
The right-hand evaluation holds above a degree threshold independent
of the map, the test subcollection, and the number of links.}
\label{fig:two-evaluations}
\end{figure}
\FloatBarrier

\subsection{Consequences for K3 geometry and zero-cycles}

The modern associated-K3 rationality prediction, in the explicit
formulation recorded by Huybrechts
\cite[Conjecture~2.1, equation~(2.1)]{HuybrechtsUpdate25}, asserts that a
smooth complex cubic fourfold is rational if and only if the primitive
cohomology of some polarized K3 surface admits a primitive isometric Hodge
embedding into its primitive fourth cohomology, with the Tate twist and
pairing convention used below.

\begin{corollary}[Failure of associated-K3 sufficiency]\label{cor:associated-k3}
Let \(d_0\) be the threshold in Theorem~\ref{thm:main}. For every admissible
\(d>d_0\), a very general \(X\in\mathcal C_d\) is irrational and admits a
labelling \(K\) of discriminant \(d\) and a primitively polarized smooth
projective K3 surface \((S_{\mathrm{Hdg}},L)\), with \(L^2=d\), for which
there is an integral Hodge isometry
\[
 K^\perp \simeq L^\perp(-1).
\]
Here \(L^\perp\subset H^2(S_{\mathrm{Hdg}},\Z)\), and the surface cup
pairing is reversed. Thus an associated untwisted polarized K3 surface in
this Hodge-theoretic sense is not sufficient for rationality.
\end{corollary}

The proof is given in Section~\ref{cub:consequences}.

\begin{corollary}[Failure of Hilbert-square sufficiency]\label{cor:hilbert-square}
Let \(d_0\) be the threshold in Theorem~\ref{thm:main}. For every integer
\(n\ge2\) such that \(d=2n^2+2n+2>d_0\), a very general
\(X\in\mathcal C_d\) is irrational and its Fano variety of lines \(F(X)\)
is birational to the Hilbert square \(S_{\mathrm{Hilb}}^{[2]}\) of a smooth
projective K3 surface \(S_{\mathrm{Hilb}}\). There are infinitely many such
discriminants. Thus birationality of \(F(X)\) to a K3 Hilbert square is
not sufficient for rationality of \(X\).
\end{corollary}

The proof is given in Section~\ref{cub:consequences}.

The categorical surface \(S\), the Hodge-associated surface
\(S_{\mathrm{Hdg}}\), and, in the Hilbert-square subfamily, the surface
\(S_{\mathrm{Hilb}}\) need not be identified. They witness different
structures on the same cubic.

\begin{corollary}[Integral Chow decompositions for the irrational cubics]
\label{cor:chow-diagonal}
Let \(d_0\) be the threshold in Theorem~\ref{thm:main}. For every admissible
\(d>d_0\), a very general \(X\in\mathcal C_d\) is irrational and has
universally trivial \(\operatorname{CH}_0\): for every field extension
\(F/\C\), the degree map
\[
 \deg\colon\operatorname{CH}_0(X_F)\longrightarrow\Z
\]
is an isomorphism. This universal \(\operatorname{CH}_0\) property is
equivalent to the existence of a point \(x\in X(\C)\), a proper closed
algebraic subset \(D\subsetneq X\), and a codimension-four cycle \(\Gamma\)
with integer coefficients supported on \(D\times X\) such that
\[
 [\Delta_X]=[X\times\{x\}]+[\Gamma]
 \quad\text{in }\operatorname{CH}^4(X\times X)
\]
with integral coefficients.
\end{corollary}

The proof is given in Section~\ref{cub:consequences}.

This applies in particular to the associated-K3 examples in
Corollary~\ref{cor:associated-k3} and the Hilbert-square subfamily in
Corollary~\ref{cor:hilbert-square}. The asserted identity is an integral
Chow-theoretic decomposition, not merely a rational cohomological
decomposition; thus the decomposition-of-the-diagonal obstruction does not
detect the irrationality of these examples. A variety is stably rational
if its product with some projective space is rational. Stable rationality
implies universal \(\operatorname{CH}_0\)-triviality \cite{Voisin17}; the
corollary supplies this necessary condition for the irrational cubics.
No assertion about stable rationality is made.

\subsection{Relation to birational invariants and quantum methods}

The existence of an associated K3 is distinct from the stronger condition
that the Fano variety of lines be birational to a Hilbert square of a K3
surface \cite[Theorems~1 and~2]{Addington16}.
A related motivic approach uses the identity of Galkin and Shinder
between a cubic, its Hilbert square and its Fano variety of lines in the
Grothendieck ring \cite[Theorem~5.1]{GalkinShinder14}.
Their rationality consequence assumes a cancellation conjecture for the
class of the affine line \cite[Conjecture~2.7 and Theorem~7.5]{GalkinShinder14};
Borisov proved that this global cancellation conjecture fails
\cite[Theorem~2.12]{Borisov18}.
Our argument uses an additive invariant of individual birational maps.

Quantum cohomology and Hodge theory provide another approach to
irrationality. Iritani's blowup theorem decomposes the quantum
\(D\)-module after a formal extension of coefficients and change of
variables \cite[Theorem~1.1]{Iritani25}. Using this quantum blowup
behavior, Katzarkov, Kontsevich, Pantev and Yu construct Hodge atoms
and obtain an irrationality theorem for a very general cubic in the full
moduli space \cite[Theorem~6.8]{KKPY26}. Their very-general hypothesis
does not cover a prescribed Hassett divisor. Gu\'er\'e obtains a
rational Hodge-theoretic constraint on a rational cubic
\cite[Theorem~63]{Guere26}; that conclusion does not assert an integral
isometry of the intersection lattices.
The argument here retains whole integral polarized transcendental
lattices through cancellation. Their integral discriminants determine
the exact degrees of the K3 surfaces in the fixed test collection.
The integral comparison and the uniform bound for the relevant surface
contributions are proved in Sections~\ref{gw:section}--\ref{sf:section}.

Integral surface-transcendental lattices also enter Kulikov's approach,
which derives irrationality of a very general cubic fourfold from an
indecomposability conjecture for those lattices
\cite[Proposition~1]{Kulikov08}. Auel, B\"ohning and Graf von Bothmer
disproved that conjecture for the sextic Fermat surface
\cite[Theorem~1.2]{AuelBohningBothmer13}.
Here the independent factors of the operator algebra distinguish whole
surface contributions even when their transcendental Hodge structures
are reducible; no general indecomposability assertion is used.

Fay proves that a very general member of a special divisor is irrational
when the quaternion class \((d/2,-3)\) over \(\Q\) is nonzero
\cite[Theorem~1.1]{Fay26}.
The class vanishes exactly when the valuation of \(d/2\) is even
at every prime congruent to two modulo three
\cite[condition~\((**')\)]{Fay26}.
For every admissible \(d\), all these valuations are zero: \(4\nmid d\)
makes \(d/2\) odd, and admissibility excludes the other such primes.
Thus Fay's quaternion obstruction vanishes throughout the range of
Theorem~\ref{thm:main}.

The divisorial invariant is a specialization of the motivic invariants of
birational maps studied by Lin and Shinder \cite{LinShinder24}.
Their later work develops horizontal and vertical versions
\cite[Definition~2.3 and Lemmas~2.4--2.5]{LinShinder26} and constructs
unbounded contributions from elliptic surfaces
\cite[Theorem~1.2 and Section~4]{LinShinder26}.
The surface calculation adapts the virtual
N\'eron--Severi viewpoint of Lin, Shinder and Zimmermann
\cite[Definition~5.2 and Proposition~5.5]{LinShinderZimmermann23}.
Our uniform exclusion concerns Picard-number-one K3 surfaces of large
degree and remains uniform when their integral transcendental
Hodge-isometry class is prescribed.

\subsection{Notation and organization}

All varieties and function fields are over \(\C\), except for the
explicit generic-field arguments. Cohomology coefficients are
specified where used; transcendental lattices always mean integral
lattices modulo torsion. For reference, the main notation is collected below.

\begin{center}
\small
\begin{tabular}{@{}p{.18\textwidth}p{.72\textwidth}@{}}
\toprule
Symbol & Meaning \\
\midrule
\(T_M,\ q_M\) & Integral transcendental middle lattice of a fourfold and its cup form. \\
\(T_Z(-1)\) & Surface transcendental lattice shifted to weight four, with pairing \(-q_Z\). \\
\(G_M\) & Rational operator algebra on \(T_M\otimes\Q\), defined from genus-zero pairings in Definition~\ref{gw:algebra}. \\
\(\mathcal S,\ u\) & A degree-\(d\) K3 test collection and its MRC indicator. \\
\(c_u(f)\) & Target-acquired minus source-lost prime divisors, weighted by \(u\). \\
\(V/B,\ Q\) & A Mori fibre-space node and the common base of a link. \\
\(v,\ h_Q\) & Integer vertical-divisor correction and its portion dominating \(Q\); \eqref{sl:v-definition} and \eqref{sl:h-definition}. \\
\(w,\ \sigma_k,\ s\) & Integer corrections on a generic surface, a conic with a chosen surface subfield, and an entire Mori node; Definition~\ref{sl:w-def}, \eqref{sf:candidate}, and \eqref{sf:node-function}. \\
\bottomrule
\end{tabular}
\end{center}

Section~\ref{pre:section} defines the test and the divisorial cocycle.
Section~\ref{gw:section} proves the whole-lattice calculation in smooth
factorizations. Section~\ref{bd:section} establishes the bounded marked
diagrams and fixes the order of the uniform choices.
Section~\ref{sl:section} computes the correction on generic surfaces;
Section~\ref{sf:section} makes it intrinsic and proves vanishing by
telescoping. Section~\ref{cub:section} applies the resulting irrationality
criterion to every sufficiently large admissible Hassett divisor and
proves the three corollaries stated above.
